\section{Discussione e Conclusioni}

\subsection{Sintesi dei Risultati Principali}
Questa ricerca ha esplorato l'efficacia di diversi LLM nella generazione e replicazione di tratti di personalità basati sul modello dei Big Five.
I risultati ottenuti offrono interessanti spunti sulle capacità e sui limiti degli LLM attuali in questo contesto specifico.

\subsubsection{Prestazioni Generali dei Modelli}
I modelli \textbf{Gemma}, in particolare i 27B, hanno dimostrato prestazioni superiori in termini di accuratezza e una notevole capacità di catturare e riprodurre le sfumature associate ai diversi tratti del modello Big Five\\ 
D'altra parte, i modelli \textbf{LLaMA}, sebbene generalmente meno accurati, hanno dimostrato una velocità di inferenza significativamente superiore, suggerendo un potenziale vantaggio in applicazioni che richiedono risposte in tempo reale.

\subsubsection{Differenze tra i Tratti di Personalità}

È emerso che alcuni tratti di personalità sono più facilmente replicabili rispetto ad altri.
Il \textbf{Nevroticismo} è risultato essere il tratto più facilmente replicabile, con correlazioni molto elevate nella maggior parte dei modelli.
Ciò indica che le \textbf{manifestazioni linguistiche} legate al Nevroticismo potrebbero essere \textbf{più distintive e riconoscibili} dai modelli di linguaggio.\\
Al contrario, l'\textbf{Estroversione} si è rivelata il tratto più difficile da replicare in modo coerente in tutti i modelli. Questa complessità potrebbe derivare dalla \textbf{natura sfaccettata del tratto}, che si esprime in modi diversi nel linguaggio scritto.

\subsubsection{Impatto della Quantizzazione}
La quantizzazione a \textbf{3 bit} (\textbf{Q3\_K\_S}) nel modello \textbf{Gemma-2-27b} ha offerto le migliori prestazioni in termini di accuratezza, ma a costo di una maggiore dimensione del modello e di un conseguente aumento del carico computazionale.\\ La quantizzazione a \textbf{2 bit} (\textbf{IQ2\_M}) ha mostrato un leggero degrado delle prestazioni, mantenendo comunque correlazioni forti per la maggior parte dei tratti, suggerendo un buon compromesso tra accuratezza ed efficienza computazionale.
\subsection{Implicazioni Teoriche e Pratiche}

\subsubsection{Comprensione della Personalità negli LLM}
La variabilità nelle prestazioni tra diversi tratti suggerisce che gli LLM potrebbero non avere una comprensione uniforme delle caratteristiche della personalità, suggerendo la presenza di potenziali bias nei dati di addestramento dei modelli.

\subsubsection{Applicazioni Pratiche}
\begin{itemize}
    \item \textbf{Sviluppo di Videogiochi}: La capacità di generare NPC con personalità coerenti e realistiche potrebbe rivoluzionare l'industria dei videogiochi, offrendo esperienze di gioco più immersive e dinamiche.
    
    \item \textbf{Chatbot e Assistenti Virtuali}: I risultati suggeriscono la possibilità di creare assistenti virtuali con personalità più ricche e coerenti, migliorando l'interazione uomo-macchina in vari contesti.
    
    \item \textbf{Ricerca Psicologica}: Gli LLM potrebbero essere utilizzati come strumenti per esplorare le manifestazioni linguistiche dei tratti di personalità, offrendo nuove prospettive per la ricerca psicologica.
\end{itemize}

\subsection{Limiti dello Studio e Direzioni Future}

\subsubsection{Limiti Metodologici}
\begin{itemize}
    \item \textbf{Dimensione del Campione}: Sebbene il numero di test effettuati sia significativo, un campione più ampio potrebbe fornire risultati ancora più robusti e generalizzabili.
    
    \item \textbf{Variabilità del Contesto}: Lo studio si è concentrato sul \textit{"50-item IPIP version of the Big Five Markers"}.
    Esplorare test più complessi o altri modelli di personalità potrebbe offrire una visione più completa delle capacità degli LLM in questo ambito.
    
    \item \textbf{Limitazioni Hardware}: La velocità di inferenza dei modelli sono state influenzate dalle limitazioni hardware, in particolare per quanto riguarda il supporto ROCm su Windows per GPU AMD.
\end{itemize}

\subsubsection{Direzioni Future}
\begin{itemize}
    \item \textbf{Esplorazione di Altri Modelli di Personalità}: Estendere lo studio ad altri modelli di personalità oltre i Big Five potrebbe offrire una comprensione più completa delle capacità degli LLM in questo ambito.
    
    \item \textbf{Analisi Linguistica Approfondita}: Un'analisi più dettagliata delle risposte generate potrebbe rivelare pattern linguistici e/o bias specifici associati a diversi tratti di personalità.
    
    \item \textbf{Fine-tuning Mirato}: Esplorare l'impatto del fine-tuning mirato su dataset specifici per la personalità potrebbe migliorare ulteriormente le prestazioni dei modelli nell'emulazione dei tratti di personalità e nella capacità "attoriale" di un modello di linguaggio.
\end{itemize}
